%!TeX root=MemoriaTFG.tex

\chapter{Ús d'eines d'intel·ligència artificial}
\label{ann:ia}

Avui en dia, l'ús d'eines d'intel·ligència artificial s'ha convertit en una pràctica
habitual tant en l'àmbit professional com en l'acadèmic. No obstant, la transparència
sobre com i per a quins propòsits s'han emprat és essencial per garantir la integritat
i la responsabilitat intel·lectual del treball. Les decisions metodològiques, les
anàlisis i les conclusions d'aquest TFG recauen íntegrament en l'autor; les eines
d'IA han actuat únicament com a suport en tasques concretes. Per aquest motiu,
a continuació s'exposa detalladament quines eines s'han emprat, per a quines fases
del TFG i amb quina finalitat.

\renewcommand{\arraystretch}{1.5}
\noindent
\begin{tabular}{@{} l @{\hspace{0.4cm}} l @{\hspace{0.4cm}} l @{}}
  \toprule
  \textbf{Eina d'IA} & \textbf{Fase del TFG} & \textbf{Ús} \\
  \midrule
  \parbox[t]{3.4cm}{Gemini (Google)\\i NotebookLM} &
  \parbox[t]{4cm}{Revisió bibliogràfica\\i estat de l'art} &
  \parbox[t]{7.2cm}{Recerca d'informació sobre algorismes d'aprenentatge per reforç
  (\acs{PPO}, \acs{DQN}, \acs{NFSP}) i marc teòric dels jocs d'informació imperfecta.} \\[30pt]

  \parbox[t]{3.4cm}{Claude Code\\(Anthropic)} &
  \parbox[t]{4cm}{Implementació i\\entrenament dels agents} &
  \parbox[t]{7.2cm}{Generació de la base dels bucles d'entrenament (\emph{training loops}).
  Els hiperparàmetres, les condicions de canvi de fase curricular i tota la lògica
  específica del joc han estat ajustats i ampliats posteriorment pel propi autor.} \\[42pt]

  \parbox[t]{3.4cm}{Claude Code\\(Anthropic)\\amb MCP Jupyter} &
  \parbox[t]{4cm}{Anàlisi i visualització\\de resultats} &
  \parbox[t]{7.2cm}{Elaboració del codi Python per generar les gràfiques dels notebooks
  d'anàlisi (corbes d'aprenentatge, mapes de calor, diagrames de radar, etc.)
  via servidor MCP de Jupyter.} \\[30pt]

  \parbox[t]{3.4cm}{GitHub Copilot\\(Microsoft)} &
  \parbox[t]{4cm}{Redacció de la memòria} &
  \parbox[t]{7.2cm}{Suport en la redacció i la revisió de textos des de l'entorn
  VS~Code. Les idees, conclusions i arguments principals han estat formulats i
  supervisats per l'autor en tot moment.} \\[4pt]
  \bottomrule
\end{tabular}
\renewcommand{\arraystretch}{1}

